\section{Localization by separated tuples}
\label{l:sep:section}

This argument measures concentration after an additional fresh block
of exact observations.  On an actual block, an auxiliary revelation
has one of two effects: it places the signal in a smaller ball, or
it produces several well-separated signals compatible with the same
data.  In the second case an exact comparison on the equal-label
constraints divides the cost of the learner's finite message by the
number of signals.  The radius decrease pays for the revelations
needed to reach that case.

Let \(S\) be uniform on \(S^{d-1}\).  In this section \(\sigma\)
denotes unnormalized surface area, and
\[
 k=\ell=\lfloor d/16\rfloor,\qquad q=\ell-1,\qquad n=d-1-k.
\]
We take \(d\) sufficiently large that \(\ell\ge2\).  A block is
\(B=(X_B,Y_B)\), where \(X_B\) has \(k\) independent standard
Gaussian rows and \(Y_B=X_BS\).  A block is fresh relative to a
variable \(U\) when its matrix is independent of the pair \((S,U)\).
We use \(h_j\) for entropy relative to unnormalized area
on a \(j\)-dimensional sphere, including spheres in affine subspaces,
and \(h\) for Lebesgue entropy.  Conditional entropies are averaged.
All logarithms are natural.

\subsection{Spherical slicing and a concentration potential}

For a full-row-rank matrix \(X\), let \(V=\operatorname{row}(X)\) and
define
\[
 R_X(s)=\sqrt{1-\|P_Vs\|^2},\qquad
 j_X(s)=\sqrt{\det(XX^{\mathsf T})}\,R_X(s).
\]
At \(R_X(s)>0\), the fiber \(\{z\in S^{d-1}:Xz=Xs\}\) is an
\(n\)-sphere of radius \(R_X(s)\).  Write \(\sigma_{X,y}\) for its
unnormalized area measure.

The normal Jacobian in the coarea formula is \(j_X(s)\); compare
Federer's Definition~2.10 and Theorem~3.1
\cite[pp.~423 and 426--427]{Federer1959}.  We give the coordinate
calculation, including the normalization needed here.  Choose
orthonormal coordinates in \(V\) and \(V^\perp\), and write
\(s=(u,R\theta)\) with \(R=\sqrt{1-\|u\|^2}\) and
\(\theta\in S^n\).  The \(\theta\)-directions scale area by \(R^n\).
The Gram matrix in the \(u\)-directions is
\(I+uu^{\mathsf T}/R^2\), of determinant \(1/R^2\), and these two
tangent groups are orthogonal.  The surface element is consequently
\(R^{n-1}\,du\,d\sigma_n(\theta)\).  The map from \(u\) to \(y=Xs\)
has absolute determinant \(\sqrt{\det(XX^{\mathsf T})}\), whereas
the area element on the fiber is \(R^n\,d\sigma_n(\theta)\).
Let \(\mathcal E_X=\{y:y^{\mathsf T}(XX^{\mathsf T})^{-1}y<1\}\).
Thus for every nonnegative Borel \(g\),
\begin{equation}
 \int_{S^{d-1}}g(s)\,\sigma(ds)
 =\int_{\mathcal E_X}\int_{\{Xs=y\}}
       \frac{g(s)}{j_X(s)}\,\sigma_{X,y}(ds)\,dy.
 \label{l:sep:area-disintegration}
\end{equation}
The right side uses only positive-radius fibers.  The omitted
critical set \(R_X=0\) has zero full-sphere
area.  The nonnegative formula follows directly from the coordinates
and Tonelli, including by exhaustion when an integral is infinite.

Suppose the relevant conditional log densities are integrable and
put \(Y=XS\).  If \(p(s\mid X,U)\) is the full-sphere density and
\(p_Y(y\mid X,U)\) the label density, then the conditional slice
density is
\[
 \frac{p(s\mid X,U)}{j_X(s)\,p_Y(Xs\mid X,U)}.
\]
Taking its logarithm gives the separate entropy identity
\begin{equation}
 h_n(S\mid X,Y,U)
 =h_{d-1}(S\mid X,U)-h(Y\mid X,U)+\mathbb E\log j_X(S).
 \label{l:sep:entropy-slicing}
\end{equation}
The area identity \eqref{l:sep:area-disintegration} and the entropy
identity \eqref{l:sep:entropy-slicing} will be used at different
points below.

Here is a finiteness check for every later use.  For a fresh Gaussian
matrix, \(R_X(S)^2\) has the law \(U/(U+V)\), where
\(U\sim\chi^2_{d-k}\) and \(V\sim\chi^2_k\) are independent.
The chi-square density at zero and its exponential tail imply
\(\mathbb E|\log R_X(S)|<\infty\).  If
\(Z_j=\sum_{i=1}^j g_i^2\) with standard Gaussians \(g_i\), then
arithmetic--geometric mean and Jensen give
\[
 \log j+\mathbb E\log g_1^2
 \le\mathbb E\log Z_j\le\log j.
\]
It follows that \(\mathbb E\log R_X(S)\ge-C\).  Gram--Schmidt on
the rows of \(X\) expresses \(\det(XX^{\mathsf T})\) as a product
of chi-squares in dimensions \(d,d-1,\ldots,d-k+1\).  Hence
\[
 \mathbb E|\log\det(XX^{\mathsf T})|<\infty,\qquad
 \mathbb E\log\det(XX^{\mathsf T})\ge k\log d-Cd.
\]
Given an unconditioned block, \(S\) is uniform on its fiber.  Its
slice area is \(|S^n|R_X(S)^n\), and its label density is that area
divided by \(j_X(S)|S^{d-1}|\).  The sampled log densities are
therefore integrable by the preceding estimates.

Adding a countable variable \(U\) with finite Shannon entropy
preserves this property.  Indeed, the average divergence of a law
conditioned further on \(U\) from the preceding conditional law is
at most \(H(U)\).  For a likelihood ratio \(v\), its negative log
part has numerator expectation at most
\(\int v(\log v)_-\,dQ\le1/e\).  The likelihood log is consequently
integrable, and so are the new log densities.  This reasoning also
works after conditioning on \(X\), even if \(U\) depends on \(X\).
We establish finite Shannon entropy for every auxiliary string before
using \eqref{l:sep:entropy-slicing} with that string.

Let \(K\) be a countable transcript with \(H(K)<\infty\).  Suppose
\(K\) specifies a closed ball \(B(z,r)\) containing \(S\), where
\(0<r\le1\), and put
\[
 \Lambda=\mathbb E[-\log r]<\infty.
\]
For an absolute \(C_v\), every affine \(n\)-sphere \(\Sigma\) obeys
\begin{equation}
 \operatorname{area}_n(\Sigma\cap B(z,r))
       \le |S^n|(C_vr)^n.
 \label{l:sep:area-bound}
\end{equation}
If the radius \(R\) of \(\Sigma\) is at most \(4r\), use its entire
area.  Otherwise, choose a point of the intersection.  The
intersection lies in its chordal cap of radius \(2r\), whose polar
angle is at most \(Cr/R\).  Polar integration bounds that area by
\(|S^{n-1}|R^n(Cr/R)^n/n\).  The sphere-area formula, or the same
polar integral, gives \(|S^{n-1}|/n\le C|S^n|\); increasing \(C_v\)
proves \eqref{l:sep:area-bound}.

For a fresh block \(B\), define
\begin{equation}
 b_n=\log|S^n|+n\log C_v,\qquad
 P=b_n-n\Lambda-h_n(S\mid K,B).
 \label{l:sep:potential}
\end{equation}
The block \(B\) is an auxiliary probe used only to evaluate this
potential.  After the transcript changes, the updated potential is
evaluated using a new independent fresh block.
The entropy of a probability density supported on a set of area
\(v\) is at most \(\log v\).  Applying this fact on each fiber and
using \eqref{l:sep:area-bound} gives \(P\ge0\).  The same conclusion
holds with a nonfresh slicing block if the conditioning still
specifies the containing ball.  With \(r=1\) and a transcript
independent of the signal and block,
\[
 P_0=n\log C_v-n\mathbb E\log R_X(S)=O(d).
\]

\begin{lemma}[Information controlled by the sliced potential]
\label{l:sep:information}
For the preceding \(K,r\) and a fresh block,
\begin{align}
 I(Y_B;K\mid X_B)&\le\frac{k}{d-1}I(S;K)+Cd,
 &h(Y_B\mid X_B)&\ge-Cd,
 \label{l:sep:subset-bounds}\\
 I(S;K)&\le C(P+d\Lambda+d),
 &-h(Y_B\mid X_B,K)&\le C(P+d\Lambda+d).
 \label{l:sep:information-bounds}
\end{align}
\end{lemma}
\begin{proof}
The subset averaging below is the entropy-cover mechanism used in
the Product Theorem of Chung, Graham, Frankl, and Shearer
\cite[Section~V, Equation~(22), pp.~33--34]{ChungGrahamFranklShearer1986}.
We give the conditional differential-entropy form required here.
Let \(Q\) be an independent Haar orthogonal frame and set
\(Z_i=\langle Q_i,S\rangle\).  Put \(N=d-1\).  For a
\(k\)-subset \(E\subset\{1,\ldots,N\}\), the entropy chain rule and
conditioning give
\[
 h(Z_E\mid Q,K)
 \ge\sum_{i\in E}h(Z_i\mid Q,K,Z_1,\ldots,Z_{i-1}).
\]
Every index occurs in the same number of \(k\)-subsets.  Averaging
over the subsets and over \(Q\), whose law is unchanged by column
permutations, yields
\begin{equation}
 h(Z_1,\ldots,Z_k\mid Q,K)
 \ge\frac{k}{N}h(Z_1,\ldots,Z_N\mid Q,K).
 \label{l:sep:subset-entropy}
\end{equation}
This uses no rotation invariance of the conditional law given \(K\).
All projected entropies are finite: the uniform projection density
has an integrable logarithm for every number of coordinates at most
\(d-1\), and finite-entropy conditioning preserves this as above.

Let \(h_j^0\) be the entropy of \(j\) fixed orthogonal coordinates
of a uniform sphere point.  The Gaussian maximum-entropy bound and
the covariance \(I_j/d\) give
\[
 h_k^0\le\frac{k}{2}\log(2\pi e/d).
\]
For \(j=N\), the two hemispherical graphs give
\[
 h_N^0=\log(|S^{d-1}|/2)+\mathbb E\log|S_d|
       \ge-\frac{N}{2}\log d-Cd.
\]
For the last inequality, use the sphere-area formula and Stirling's
bound, and write \(S=g/\|g\|\); Jensen bounds
\(\mathbb E\log\|g\|\) above by \(\tfrac12\log d\), whereas
\(\mathbb E\log|g_d|\) is a finite absolute constant.
Data processing gives
\[
 h(Z_1,\ldots,Z_N\mid Q,K)\ge h_N^0-I(S;K).
\]
Thus \eqref{l:sep:subset-entropy} bounds the information in \(k\)
frame coordinates by \(\tfrac{k}{N}I(S;K)+Cd\).

Given \(X_B\), its label is an invertible linear transform of the
orthogonal coordinates in its row space.  That row space is uniform,
independent of \((S,K)\); extending an orthonormal row frame gives
the \(Q\) experiment above.  Hence the same information bound holds
for \(I(Y_B;K\mid X_B)\).  Without \(K\),
\eqref{l:sep:subset-entropy} gives
\(h_k^0\ge(k/N)h_N^0\).  Passing from orthogonal coordinates to
\(Y_B\) adds
\(\tfrac12\mathbb E\log\det(X_BX_B^{\mathsf T})\), which cancels
the leading \(-\tfrac{k}{2}\log d\) term.  This proves both
inequalities in \eqref{l:sep:subset-bounds}.

By the definition of \(P\),
\[
 I(S;K\mid B)=P-P_0+n\Lambda.
\]
Freshness and \(Y_B=X_BS\) give the chain-rule identity
\[
 I(S;K)=I(S;K\mid B)+I(Y_B;K\mid X_B).
\]
Since \(k/(d-1)\le1/4\), the first inequality of
\eqref{l:sep:subset-bounds} can be absorbed on the left.
This proves the information estimate in
\eqref{l:sep:information-bounds}.  Finally
\(h(Y_B\mid X_B,K)=h(Y_B\mid X_B)-I(Y_B;K\mid X_B)\);
the two bounds just proved give the negative-entropy estimate.
\end{proof}

\subsection{A terminating revelation and a separated tuple}

Let \(A=(X_A,Y_A)\) be the actual next block, fresh relative to \(K\).
Let \(W\) be its outgoing message, with at most \(e^b\) values.
The transcript includes the incoming learner state, so the law of
\(W\) is a kernel of \((K,A)\).  The next construction uses its own
auxiliary randomness.  It does not change the learner's rule.

For a tuple \(\mathbf s=(s_1,\ldots,s_\ell)\), define
\[
 D_{\mathbf s}=[s_2-s_1,\ldots,s_\ell-s_1],\qquad
 \mathsf V(\mathbf s)=\sqrt{\det(D_{\mathbf s}^{\mathsf T}D_{\mathbf s})}.
\]
This is \(q!\) times the simplex volume and is symmetric in its
points.  Gram--Schmidt expresses it as the product of successive
distances to predecessor affine spans.

\begin{lemma}[Entropy cost of obtaining a separated tuple]
\label{l:sep:localization}
There are fixed absolute constants \(a\ge1\), \(\delta>0\), and
\(b_0>0\) with the following property.  Put
\(\eta=e^{-a}\), \(p_0=1/(4\ell)\), and
\(u_d=\log(1/p_0)+1\).  There is a countable string \(C_*\), generated
from \((K,A,S)\) and separate randomness, and an integrable integer
\(N\ge0\), such that \(K,C_*\) specifies a ball of radius
\(r_*=re^{-aN}\) containing \(S\).  The string is conditionally
independent of \(W\) given \((K,A,S)\), and in fact
\(W\) is independent of \((S,C_*)\) given \((K,A)\).  Moreover,
\begin{align}
 H(C_*\mid K,A)&\le(b_0d+\ell a+u_d)\mathbb E N+u_d,
 \label{l:sep:conditional-code}\\
 H(C_*\mid K)&\le C_a d(1+\mathbb E N),
 \label{l:sep:ambient-code}\\
 0\le P_{\mathrm{mid}}
 &:=b_n+n\mathbb E\log r_*-h_n(S\mid K,A,C_*) \notag\\
 &\le P-\delta da\,\mathbb E N+u_d.
 \label{l:sep:refinement-payment}
\end{align}
Conditional on \((K,A,C_*)\), there is a symmetric law \(Q\) of
\(\ell\) points, each with the conditional signal marginal, such that
\begin{equation}
 \mathsf V(\mathbf s)\ge(\eta r_*)^q,\qquad
 D\!\left(Q\middle\|
       \mathcal L(S\mid K,A,C_*)^{\otimes\ell}\right)\le\log2.
 \label{l:sep:passing-tuple}
\end{equation}
\end{lemma}
\begin{proof}
We describe a decision tree conditional on \(K,A\).  At a node let
\(\nu\) be the signal law after the recorded prefix, supported in
the current ball \(B(z_0,r_0)\).  Enumerate once and for all a
countable dense family of affine subspaces of dimensions at most
\(\ell-2\), for example those generated by rational points.
If the first subspace \(E\) with
\[
 \nu\{s:\operatorname{dist}(s,E)<2\eta r_0\}>p_0
\]
exists, record whether the signal belongs to this tube.  On a negative
answer, continue with the conditional law and the same ball.  On a
positive answer, cover the current ball by a fixed ordered maximal
\(\eta r_0\)-separated net, choose the first center within
\(\eta r_0\) of \(S\), record its index, and replace the ball by the
radius-\(\eta r_0\) ball about that center.  Then continue.

The full ambient net has at most \((1+2/\eta)^d\) points.  Given
the selected subspace, only \(\exp(b_0d+\ell a)\) indices can occur,
where \(b_0\) is independent of \(a\ge1\).  To see this, an eligible
center lies within \(3\eta r_0\) of \(E\).  Its projection to \(E\)
lies in a radius-\(r_0\) ball, coverable by
\((1+2/\eta)^v\) radius-\(\eta r_0\) balls when \(v=\dim E\).
Over each covering ball, the eligible centers lie in an ambient
radius-\(4\eta r_0\) ball.  Their disjoint radius-\(\eta r_0/2\)
balls fit in its radius-\(9\eta r_0/2\) enlargement, giving at most
\(9^d\) centers.  Since \(v\le\ell-2\), the claimed exponential
bound follows after enlarging \(b_0\).

If no selected subspace exists, every affine subspace of dimension
at most \(\ell-2\) has mass at most \(p_0\) in its open
\(\eta r_0\)-tube.  Indeed its distance function on the bounded
current ball can be approximated uniformly within \(\eta r_0\)
by that of a subspace from the countable family.  Mass greater than
\(p_0\) in the smaller open tube would therefore trigger the larger
tube test for a member of the family.

In this case, draw \(\ell-1\) additional independent points from
\(\nu\), using \(S\) as the first point, and record whether
\[
 \mathsf V(s_1,\ldots,s_\ell)\ge(\eta r_0)^q.
\]
On failure continue at the same scale with the updated signal law;
on success stop.  Before the outcome, the tuple is i.i.d. from
\(\nu\).  If the volume is smaller than \((\eta r_0)^q\), one of
the \(q\) successive distances is smaller than \(\eta r_0\).
Conditionally on its predecessors, each such event has probability
at most \(p_0\).  The pass probability is therefore at least
\(1-qp_0\ge1/2\).

Record the decision type, its outcome, and any ambient net index,
but not the extra test points.  Given \((K,A)\) and the prefix, the
type and selected subspace are already determined.  The first
successful subspace is a measurable choice because the family is
countable and its tube probabilities are kernel integrals of Borel
indicators.  Fixed scaled nets and ordered ties give measurable
centers.  Regular conditional laws at nodes, with arbitrary defaults
at null nodes, complete a recursion by probability kernels on
standard Borel spaces.  Independent auxiliary draws make
\(W\) independent of \((S,C_*)\) given \((K,A)\).

First truncate the tree after \(u\) decisions, and write \(C_u,N_u\)
for its string and number of refinements.  At every active node the
progress outcome, meaning refinement or a successful test, has
conditional probability \(p\ge p_0\).  Its binary entropy is at most
\[
 p\{\log(1/p_0)+1\}=p u_d,
\]
because \(-p\log p\le p\log(1/p_0)\) and
\(-(1-p)\log(1-p)\le p\).  The expected number of progress outcomes
is at most \(\mathbb E N_u+1\).  Summing the node entropies and the
eligible index costs gives
\begin{equation}
 H(C_u\mid K,A)
 \le(b_0d+\ell a+u_d)\mathbb E N_u+u_d.
 \label{l:sep:truncated-code}
\end{equation}
The expected number of active decisions is at most
\(p_0^{-1}(\mathbb E N_u+1)\).  For an encoding that is decoded
given \(K\) alone, use two bits for the type and outcome at each
active node and a fixed-length full ambient net index at a refinement.
The decoder knows the truncation depth \(u\); the recorded type and
outcome determine any earlier stop.  The ambient index recovers the
next center without knowing the selected subspace.
This prefix encoding has expected length at most
\(C_a d(1+\mathbb E N_u)\) nats, so
\begin{equation}
 H(C_u\mid K)\le C_a d(1+\mathbb E N_u).
 \label{l:sep:truncated-ambient}
\end{equation}

At finite depth all strings have finite entropy.  The area bound and
the conditional mutual-information identity give
\begin{align*}
 0\le P_{\mathrm{mid},u}
 &=P-na\mathbb E N_u+I(S;C_u\mid K,A)\\
 &\le P-\{(n-\ell)a-b_0d-u_d\}\mathbb E N_u+u_d.
\end{align*}
For large \(d\), \(n-\ell>d/2\) and \(u_d=O(\log d)\).
Choose \(a\) once and for all large enough that the coefficient in
braces is at least \(\delta da\) with an absolute \(\delta>0\).
It follows that \(\sup_u\mathbb E N_u<\infty\), and then that the
expected number of active decisions is bounded uniformly in \(u\).
By monotone convergence the full number of decisions is integrable,
so the tree terminates almost surely.  Denote its string and
refinement count by \(C_*,N\).

The same prefix encoding at termination has finite expected length,
so \(H(C_*\mid K)<\infty\).  The truncated prefixes increase to the
full string and \(N_u\uparrow N\).  Conditional information in the
prefixes increases to the information in the full countable string;
this also follows from the chain rule and
\(H(C_*\mid C_u,K,A)\to0\), using its finite entropy.
Passing to the limit in the preceding bounds proves
\eqref{l:sep:conditional-code}--\eqref{l:sep:refinement-payment}.
Thus the stopping argument did not presume entropy finiteness for
an infinite procedure.

Finally condition on a terminal prefix just before its successful
test, with signal law \(\nu\).  Let \(E_{\mathrm{pass}}\) be its
symmetric volume event and set
\(Q=\nu^{\otimes\ell}(\,\cdot\mid E_{\mathrm{pass}})\).
The symmetry makes all marginals equal, and the first marginal is
exactly the signal law after recording the pass.  Call it \(\mu_*\).
If \(p=\nu^{\otimes\ell}(E_{\mathrm{pass}})\ge1/2\), then the
product-reference decomposition gives
\[
 D(Q\|\mu_*^{\otimes\ell})
 =D(Q\|\nu^{\otimes\ell})-\ell D(\mu_*\|\nu)
 \le\log(1/p)\le\log2.
\]
All terms are finite because \(dQ/d\nu^{\otimes\ell}
=\mathbf1_{E_{\mathrm{pass}}}/p\).  The volume bound holds throughout
the conditioned law.  This proves \eqref{l:sep:passing-tuple}.
\end{proof}

\subsection{An exact comparison on the equal-label constraints}

Use the coupling of Lemma~\ref{l:sep:localization} conditional on
\((K,A,C_*)\), and draw a common \(W\) from the learner's kernel
given \((K,A)\).  Every \((S_i,K,A,C_*,W)\) has the true signal
marginal.  Write
\[
 h_f=h_{d-1}(S\mid K,C_*),\qquad
 h_{\mathrm{mid}}=h_n(S\mid K,A,C_*),
\]
and let \(D_{\mathrm{coup}}\le\log2\) be the averaged divergence
in \eqref{l:sep:passing-tuple}.  We call divergence of a tuple law
from the product of its marginals total correlation, following
Watanabe~\cite{Watanabe1960}; the probability-law formulation
and product decompositions are given by
\cite[Section~5, Equation~(26) and Proposition~5.1]{Austin2020}.

Conditional on \((K,C_*)\), form a reference law by drawing
\(\ell\) independent points from \(\mathcal L(S\mid K,C_*)\), then
drawing each row of \(X\) as a standard Gaussian in
\((\operatorname{col}D_{\mathbf s})^\perp\).  Let \(\mathcal E\)
be the averaged relative entropy of the actual coupled tuple and
matrix \(X_A\) from this reference.

\begin{lemma}[The tuple comparison]
\label{l:sep:comparison}
The divergence \(\mathcal E\) is finite and
\begin{align}
 \mathcal E
 &=I(X_A;C_*\mid S,K)
   +q\{h_f-h_{\mathrm{mid}}+\mathbb E\log j_{X_A}(S)\}
       \notag\\
 &\quad-\frac{kq}{2}\log(2\pi)
       -k\mathbb E\log\mathsf V(\mathbf S)+D_{\mathrm{coup}}.
 \label{l:sep:comparison-identity}
\end{align}
Furthermore,
\begin{equation}
 \ell I(S;W\mid K,C_*)\le b+\mathcal E.
 \label{l:sep:message-sharing}
\end{equation}
\end{lemma}
\begin{proof}
We first identify a common measure for the two singular laws.  Let
\(\mathcal R\) be the open region in the ambient tuple--matrix space
where \(D_{\mathbf s}\) has rank \(q\), \(X\) has rank \(k\), and
\(j_X(s_i)>0\) for every \(i\).  On the incidence constraint
\(XD_{\mathbf s}=0\), let \(\gamma_{\mathbf s}\) be the Gaussian
matrix law conditioned on that constraint.  The ambient Gaussian
density of \(XD_{\mathbf s}\) at zero is
\[
 \phi_{\mathbf s}=(2\pi)^{-kq/2}\mathsf V(\mathbf s)^{-k}.
\]
Indeed each row has covariance
\(D_{\mathbf s}^{\mathsf T}D_{\mathbf s}\).  Write \(p_X\) for
ambient Gaussian matrix density.  On the regular incidence set the
locally finite measure
\begin{equation}
 \sigma^{\otimes\ell}(d\mathbf s)\,
       \phi_{\mathbf s}\,\gamma_{\mathbf s}(dX)
 \label{l:sep:common-tuple-first}
\end{equation}
also has the integration description
\begin{equation}
 p_X(X)\,dX\,\sigma(ds_1)
       \prod_{i=2}^{\ell}
       \frac{\sigma_{X,Xs_1}(ds_i)}{j_X(s_i)}.
 \label{l:sep:common-matrix-first}
\end{equation}

To verify this equality, take a continuous test \(G(\mathbf s,X)\)
compactly supported in \(\mathcal R\).  In its integral under
\(\sigma^{\otimes\ell}(d\mathbf s)p_X(X)dX\), insert the normalized
indicator of a radius-\(\varepsilon\) ball in \(\mathbb R^{kq}\),
evaluated at \(XD_{\mathbf s}\).  Fixing the tuple first, the
conditional law given \(XD_{\mathbf s}=Z\) is
\(\gamma_{\mathbf s}\) translated by
\(Z(D_{\mathbf s}^{\mathsf T}D_{\mathbf s})^{-1}D_{\mathbf s}^{\mathsf T}\).
This translation and the Gaussian image density vary continuously
at \(Z=0\).  On the tuple projection of the compact support, the
smallest singular value of \(D_{\mathbf s}\) is bounded away from
zero, so the image densities are uniformly bounded.  Dominated
convergence gives \eqref{l:sep:common-tuple-first}.

In the other order, fix \(X,s_1\) and apply the area identity
\eqref{l:sep:area-disintegration} separately to
\(s_2,\ldots,s_\ell\), using labels
\(X(s_i-s_1)\).  The affine fibers near the zero label are
continuous translations and scalings of unit spheres.  On the
compact regular support their radii and \(j_X(s_i)\) are bounded
away from zero, so the fiber integrals of \(G\) are continuous at
zero and uniformly bounded.  Dominated convergence in the label
vector gives \eqref{l:sep:common-matrix-first}.  Equality on all
continuous compactly supported tests identifies the locally finite
Borel measures on the regular set.  Thus it also holds for
nonnegative measurable integration there.  This is a regular-region
identity; no global finite mass for the unweighted incidence measure
has been assumed.

Both probability laws in the comparison give this region full mass.
For the reference, the one-point law has a density relative to
sphere area, so its independent tuple is affinely independent
almost surely.  Its constrained rows live in dimension \(d-q>k\)
and have full row rank almost surely.  On the constraint, all
points have the same row-space projection.  That projection has
norm strictly less than one: if \(s_1\) were in the row space,
then \(\langle s_1,s_i-s_1\rangle=0\) would give
\(\langle s_1,s_i\rangle=1\), forcing \(s_i=s_1\) and contradicting
full affine rank.  For the actual law, the volume lower bound gives
full affine rank, and the true \((S_1,X_A)\) marginal gives full
row rank and positive residual radius almost surely.  The quantities
\(j_X(s_i)\) agree on the constraint because their row-space
projections agree.

Condition now on \((K,C_*)\).  Jointly measurable density versions
can be obtained by taking Radon--Nikodym derivatives of the joint
laws with their conditioning variables and then disintegrating;
absolute continuity follows from the finite-entropy conditioning
above and from the coupling bound.  Let \(p_c\) be the full-sphere
marginal density, \(p_{1,X}\) the density of the true
\((S_1,X_A)\) law relative to \(\sigma\otimes dX\), and
\(p_{\mathrm{sl}}(\,\cdot\mid A)\) the conditional slice density.
Let \(t(\mathbf s\mid A)\) be the coupling density relative to the
product of these slice marginals.  The reference density relative
to \eqref{l:sep:common-tuple-first} is
\[
 \frac{\prod_{i=1}^{\ell}p_c(s_i)}{\phi_{\mathbf s}}.
\]
The actual density relative to the same measure, using
\eqref{l:sep:common-matrix-first}, is
\[
 \frac{p_{1,X}(s_1,X)}{p_X(X)}
       \prod_{i=2}^{\ell}
       \{p_{\mathrm{sl}}(s_i\mid A)j_X(s_i)\}\,
       t(\mathbf s\mid A).
\]
Here \(A=(X,Xs_1)\) is determined by \(X,s_1\).  Each \(p_c(S_i)\)
is finite and positive almost surely under the actual law, since
each coordinate has that marginal after averaging \(A\).  These
densities therefore establish absolute continuity of the actual
law with respect to the reference, rather than transferring
arbitrary null sets between the singular laws.

The first-coordinate contribution to the expected log ratio is
\begin{align*}
 \mathbb E\log\frac{p_{1,X}(S_1,X_A)}
                         {p_X(X_A)p_c(S_1)}
 &=I(S;X_A\mid K,C_*)+I(X_A;C_*\mid K)\\
 &=I(X_A;C_*\mid S,K).
\end{align*}
The last equality uses independence of \(X_A\) from \((S,K)\) at
block entry.  Each of the other \(q\) coordinates contributes
\(h_f-h_{\mathrm{mid}}+\mathbb E\log j_{X_A}(S)\).
The term \(\log\phi_{\mathbf S}\) contributes
\(-\tfrac{kq}{2}\log(2\pi)-k\mathbb E\log\mathsf V\), and
\(\log t\) contributes \(D_{\mathrm{coup}}\).
These are finite log integrals.  The signal and slice log densities
are integrable because \(H(K,C_*)<\infty\); the first-coordinate
divergence is at most \(H(C_*\mid K)\); the Jacobian has its original
true marginal; and
\[
 q(\log\eta+\log r_*)\le\log\mathsf V(\mathbf S)\le q\log2
\]
with \(\mathbb E[-\log r_*]<\infty\).  The coupling log likelihood
is integrable by its finite divergence and the usual \(1/e\)
negative-part bound.  Taking the log ratio now proves
\eqref{l:sep:comparison-identity}.

Discard \(X_A\) by data processing.  The divergence of the tuple
from the product of its \((K,C_*)\)-conditional marginals is at
most \(\mathcal E\).  Adding \(W\) and using an independent
\((K,C_*)\)-conditional copy of \(W\) in that reference adds at
most \(H(W\mid K,C_*)\le b\).  Decomposing the same finite
divergence after conditioning on \(W\) gives
\[
 \ell I(S;W\mid K,C_*)+
 \mathbb E D\!\left(
    \mathcal L(\mathbf S\mid K,C_*,W)
    \middle\|
    \mathcal L(S\mid K,C_*,W)^{\otimes\ell}\right).
\]
The equality of every signal--message marginal, established before
the lemma, is what makes the first term \(\ell\) times the same
information.  The second term is nonnegative.  This proves
\eqref{l:sep:message-sharing}.
\end{proof}

\subsection{Cancellation and iteration}

Set \(K_+=(K,C_*,W)\), and let \(B\) now be fresh relative to
\(K_+\).  Write
\[
 h_+=h_n(S\mid K_+,B),\qquad
 \Lambda_+=\Lambda+a\mathbb E N.
\]
The entropy identity \eqref{l:sep:entropy-slicing} gives
\begin{align*}
 h_+&=h_f-I(S;W\mid K,C_*)-h(Y_B\mid X_B,K_+)
                         +\mathbb E\log j_{X_B}(S),\\
 h_{\mathrm{mid}}-h_f-\mathbb E\log j_{X_A}(S)
   &=-h(Y_A\mid X_A,K,C_*)-I(S;X_A\mid K,C_*).
\end{align*}
The expected Jacobians agree because, before conditioning, each
matrix has the standard Gaussian law independent of the uniform
signal.  Subtract the first identity from \(h_{\mathrm{mid}}\),
then use \eqref{l:sep:comparison-identity} and
\eqref{l:sep:message-sharing}.  The coefficient of the
mid-versus-full entropy difference is \(1-q/\ell=1/\ell\).
The chain rule and entry independence give
\[
 I(X_A;C_*\mid S,K)-I(S;X_A\mid K,C_*)
       =I(X_A;C_*\mid K).
\]
The resulting inequality is
\begin{align}
 h_{\mathrm{mid}}-h_+
 &\le h(Y_B\mid X_B,K_+)
       -\frac{h(Y_A\mid X_A,K,C_*)}{\ell}
       -\frac{kq}{2\ell}\log(2\pi) \notag\\
 &\quad-\frac{k}{\ell}\mathbb E\log\mathsf V(\mathbf S)
       +\frac{I(X_A;C_*\mid K)+b+D_{\mathrm{coup}}}{\ell}.
 \label{l:sep:entropy-comparison}
\end{align}
In particular, the selection of the actual rows has left the
explicit cost \(I(X_A;C_*\mid K)\); those rows have not been treated
as fresh after selection.

Let \(z_*\) be the decoded final center.  Compare the fresh label
law to \(N(X_Bz_*,r_*^2I_k)\).  Since \(X_B\) is independent of
\((S,K_+)\) and \(\|S-z_*\|\le r_*\),
\[
 h(Y_B\mid X_B,K_+)
 \le k\mathbb E\log r_*+\frac{k}{2}\log(2\pi e).
\]
Together with \(\log\mathsf V\ge q(\log\eta+\log r_*)\), the
scale-dependent terms in \eqref{l:sep:entropy-comparison} are at most
\[
 k\mathbb E\log r_*-\frac{kq}{\ell}\mathbb E\log r_*+O_a(k)
 =\frac{k}{\ell}\mathbb E\log r_*+O_a(k)\le O_a(k).
\]
The residual coefficient is positive and \(\log r_*\le0\).
For the actual labels,
\[
 -h(Y_A\mid X_A,K,C_*)
 \le-h(Y_A\mid X_A,K)+H(C_*\mid K),\qquad
 I(X_A;C_*\mid K)\le H(C_*\mid K).
\]
Apply Lemma~\ref{l:sep:information} and
\eqref{l:sep:ambient-code}.  Since \(\ell\ge d/32\) for large
\(d\), when \(b\le d^2\) we obtain
\begin{equation}
 h_{\mathrm{mid}}-h_+
 \le C_a d+\frac{C'}d(P+d\Lambda)+C_a\mathbb E N.
 \label{l:sep:one-block}
\end{equation}

\begin{proposition}[The controlled-horizon recurrence]
\label{l:sep:recurrence}
Suppose the outgoing message has at most \(e^b\) values, with
\(b\le d^2\).  For the localization update above, there are absolute
\(\gamma,C_1,C_2>0\) such that
\begin{equation}
 P_++\gamma d\Lambda_+
 \le(1+C_1/d)(P+\gamma d\Lambda)+C_2d.
 \label{l:sep:recurrence-formula}
\end{equation}
Starting with radius one and signal-independent initial information,
after \(m\le d\) blocks whose messages have at most \(e^{d^2}\)
values,
\begin{equation}
 P_m+\gamma d\Lambda_m\le C d(m+1),\qquad
 I(S;K_m)\le C d(m+1).
 \label{l:sep:horizon}
\end{equation}
\end{proposition}
\begin{proof}
By definition,
\(P_+=P_{\mathrm{mid}}+h_{\mathrm{mid}}-h_+\).
Combine \eqref{l:sep:refinement-payment} and
\eqref{l:sep:one-block}, and take \(\gamma=\delta/2\).
For sufficiently large \(d\), the coefficient
\(-\delta da+C_a+\gamma da\) of \(\mathbb E N\) is nonpositive.
The term \(u_d=O(\log d)\) is absorbed in \(C_2d\).  Since
\(P,\Lambda\ge0\), the remaining
\(C'(P+d\Lambda)/d\) is bounded by
\((C_1/d)(P+\gamma d\Lambda)\).  This proves
\eqref{l:sep:recurrence-formula}.

Every new transcript is countable with finite entropy, specifies the
new containing ball, and has \(\Lambda_+<\infty\).  Fresh future
matrices can therefore be generated and the construction repeated.
Starting from \(P_0=O(d)\) and \(\Lambda_0=0\), iteration through
\(m\le d\) steps has multiplier
\((1+C_1/d)^m\le e^{C_1}\), which proves the potential bound in
\eqref{l:sep:horizon}.  Lemma~\ref{l:sep:information} then gives
the information bound.  No linear-growth claim for unrestricted
numbers of blocks is made.
\end{proof}

\subsection{The full accuracy range}

\begin{theorem}[The sample bound by separated-tuple localization]
\label{l:sep:theorem}
Let \(M(d)=o(d^2)\) be integer-valued and let
\(0<\epsilon(d)\le1/10\).  Every learner of
Section~\ref{l:model} with uniform-sphere angular success probability
at least \(2/3\) must satisfy
\[
 T(d)\ge c\,d\log(1/\epsilon(d))
\]
for all sufficiently large \(d\), with an absolute \(c>0\).
The eventual dimension threshold may depend on the memory sequence.
Success at least \(2/3\) for every fixed signal implies the same
conclusion by averaging.
\end{theorem}
\begin{proof}
Set \(L=\log(1/\epsilon)\).  For almost every realization of shared
rule randomness, Lemmas~\ref{l:borel} and \ref{l:prior-ac} give
Borel rules with the same uniform-prior experiment.  All bounds
below are numerical bounds on that conditional success probability,
with constants independent of the realization, so they may be
averaged.  No jointly selected family of Borel versions is needed.
The joint independence of the shared randomness and initial state
from the signal and rows ensures the required signal-independent
initial transcript for almost every value of the shared randomness.

We first explain why the restriction \(m\le d\) suffices.  An
angular cap of radius \(\epsilon\) has uniform mass at most
\((2\epsilon)^{d-1}\).  For example, its polar-area numerator is
at most \(\epsilon^{d-1}/(d-1)\), while its denominator is at least
\((\pi/3)(\sqrt3/2)^{d-2}\); the displayed bound follows for large
\(d\).  There are at most \((T+1)2^M\) terminal time--state pairs.
For each pair, integration over its fixed fresh output law contributes
at most one cap mass, even before multiplying by its reach probability.
Summing and averaging over shared randomness gives
\begin{equation}
 (d-1)L\le M\log2+\log(T+1)+(d-1)\log2+\log(3/2).
 \label{l:sep:output-count}
\end{equation}
At indices with \(T<dL\),
\(\log(T+1)\le\log(d+1)+L\).  Since \(M=o(d^2)\),
\eqref{l:sep:output-count} implies \(L=o(d)\) along those indices.
Also
\[
 b=\log((T+2)2^M)\le d^2
\]
eventually, and after padding to a multiple of \(k\),
\[
 m=\lceil T/k\rceil\le 32L+1=o(d)\le d
\]
eventually.  Thus every potentially short run lies within the
proved horizon.

At indices with \(T\ge dL\), the claimed conclusion holds with
constant one.  At all remaining indices, the horizon and width
bounds just established apply eventually.  Pad an early-stopping
program by retaining its terminal state and stopping index, then
ignore later observations and any suffix added to complete the last
block.  At boundaries its state is a message with at most \(e^b\)
values.  The augmented transcript retains all earlier auxiliary
strings, but the learner uses only its state.  Future matrices remain
fresh, so Proposition~\ref{l:sep:recurrence} applies with
\(m=\lceil T/k\rceil\le d\).  For each shared seed,
\(I(S;K_m)\le C_0d(m+1)\), with an absolute \(C_0\).

Let \(\pi\) be the success probability for that seed.  The output
uses information contained in \(K_m\) and fresh randomness.  Under
the product of its marginal and the uniform signal law, success
has probability at most \((2\epsilon)^{d-1}\).  By data processing
and \eqref{l:binary},
\[
 I(S;K_m)\ge
 \pi(d-1)\log\frac1{2\epsilon}-\log2.
\]
Combining the two bounds gives a numerical inequality in \(\pi\);
averaging it over the shared seed replaces \(\pi\) by the original
success probability, at least \(2/3\).  Thus, for absolute
\(L_0,c_0>0\), whenever \(L\ge L_0\),
\[
 c_0dL\le C_0d(m+1)\le C_0d(T/k+2),
\]
so \(T\ge k(c_0L/C_0-2)\).  Choose
\(L_1=\max\{L_0,4C_0/c_0\}\).  Since \(k\ge d/32\), this gives
\[
 T\ge \frac{c_0}{64C_0}\,dL\qquad(L\ge L_1).
\]
For \(L<L_1\), the local full-data endpoint in
Lemma~\ref{l:caps-endpoint} gives \(T>d/4\), hence
\(T\ge dL/(4L_1)\).  The minimum of these absolute constants and
one works for every allowed accuracy sequence, including
nonmonotone ones.
\end{proof}

\begin{remark}[Vanishing success with sublinear samples]
The full-data calculation also gives a stronger small-sample
conclusion.  For \(0\le T\le d-3\), reveal all \(T\) rows and
labels.  As in Lemma~\ref{l:caps-endpoint}, the conditional signal
is uniform on a residual sphere of radius \(R\), with
\(\Pr\{R<1/2\}\le4T/(3d)\).  On \(R\ge1/2\), a radius-\(\epsilon\)
error ball meeting that sphere is contained in a chordal cap of
radius \(2\epsilon\) about a point of the intersection.  For
\(\epsilon\le1/10\), its polar angle is at most
\(\theta_0=2\arcsin(1/5)\).  Since \(\sin\theta_0<2/5\), polar
integration bounds its conditional mass by
\[
 \frac{\int_0^{\theta_0}\sin^{d-T-2}\theta\,d\theta}
      {\int_0^\pi\sin^{d-T-2}\theta\,d\theta}
 \le C\left(\frac4{5\sqrt3}\right)^{d-T-2}.
\]
Consequently every estimate based on these data and independent
randomness has success probability at most
\[
 \frac{4T}{3d}+C\left(\frac4{5\sqrt3}\right)^{d-T-2}.
\]
This tends to zero when \(T=o(d)\), uniformly for
\(0<\epsilon\le1/10\), and therefore also applies to early-stopping
finite-state learners.
\end{remark}
